% ================================================================
% gear_fatigue_v3.0.tex — Gear Bending Fatigue Variance Decomposition
% v3.0: Full rewrite — AFT numbers corrected to EM-estimated ground truth,
%       Bootstrap 2000 draws, Discussion narrative: Three Co-Equal Factors,
%       multi-torque downgraded to qualitative, We→I throughout
% v2.9: ANOVA→AFT (censored likelihood), but incomplete — numbers mixed old/new
% v2.8: Bootstrap numbers corrected; bucur→vietze; 5 reviewer fixes
% v2.7: +Bonaiti (2024) +Vietze (2024) experimental citations,
%       bootstrap 100→2000 fix (Abstract/figures only), expanded Discussion §4.3
% v2.6: Torque-dependence + repositioned as general UQ framework
% v2.5: 3 experimental papers cited (Wang, Pagliari, Glodez)
% v2.4: Kf removed, torque=700 Nm, 5 switches, 144 combos
% All formulas verified against ISO 6336-3:2019, AGMA 2001-D04, FKM 7th ed.
% S-N data from Waterloo SMDIdbase (publicly accessible)
% All numbers from output/aft_anova.json (EM estimation, 2000 bootstrap) — v3.0
% Target: International Journal of Fatigue (Elsevier)
% Versions preserved: v1.0, v2.0, v2.1, v2.2, v2.3, v2.4, v2.5, v2.6, v2.7, v2.8, v2.9
% ================================================================
\documentclass[12pt,a4paper]{article}
\usepackage[margin=2.5cm]{geometry}
\usepackage{amsmath,amssymb,graphicx,booktabs,multirow,hyperref}
\usepackage[numbers,sort&compress]{natbib}

\begin{document}

\title{\textbf{Method-Induced Uncertainty in Gear Bending Fatigue:\\
Factorial Variance Decomposition Across Five Engineering Choices}}

\author{Zhilei Chen\\
Guangdong Peizheng College, Data and Computer Science,\\
53 Peizheng Rd., Chini Town, Huadu, Guangzhou, Guangdong 510830, China\\
E-mail: 2604513@peizheng.edu.cn}

\maketitle

% ================================================================
% @sources: aft_anova.json (variance fractions, EM estimation, 2000 bootstrap), sweep_summary.json (Nf range, counts), gear_params.py (geometry, material)
\begin{abstract}
\label{sec:abstract}
% ================================================================

When three engineers compute the bending fatigue life of the same spur
gear---same geometry, same material, same load---their predictions can
differ by four orders of magnitude. The divergence comes not from
material variability or manufacturing defects, but from the
methodological choices each engineer makes: which standard to apply, how
to penalize surface roughness, which S--N curve to trust, and how to
correct for mean stress. I hold a single spur gear pair
(module~3~mm,% @source: module_mn
20 teeth,% @source: teeth_pinion
42CrMo4/AISI~4140,% @source: material
700~N$\cdot$m,% @source: torque_Nm
constant amplitude, $R=0$) fixed and systematically vary five methodological% @source: stress_ratio_R
choices: size-and-surface standard (ISO~6336, AGMA~2001, FKM), surface
roughness grade (Rz~=~4, 15, 50~$\mu$m), S--N curve source (Waterloo
SMDIdbase \#68, \#66), mean-stress correction (Goodman, Gerber, Morrow,
SWT), and cumulative damage rule (Miner, Modified Miner). The tooth
root stress concentration is handled by the ISO~6336 $Y_{Sa}$ factor
within the nominal stress calculation; a separate fatigue notch factor
$K_f$ was removed to avoid double-counting of the notch effect. A
full-factorial AFT survival model over all 144 combinations (including
28 run-outs handled via censored likelihood) decomposes the total
log-variance. Three factors share the dominant contribution: mean-stress
correction (33.9\%), size-and-surface standard (30.1\%), and surface
roughness (24.6\%). A Standard~$\times$~Rz interaction accounts for
3.0\%. The S--N curve source (7.5\%) and cumulative damage rule (0.9\%)
are secondary. Bootstrap validation confirms stable
rankings.% @source: n_bootstrap
The factor ranking is operating-point-dependent: at lower torque
the standard choice dominates, while mean-stress correction and
surface roughness rise to dominance near the endurance-limit
regime explored here. Whereas prior studies---including recent work in this
journal~\cite{bonaiti2024}---ask whether two standards produce different
predictions, this study asks how much each methodological switch contributes
to total log-variance relative to all others, by simultaneous factorial
decomposition across five dimensions. No prior work has ranked the
sources of method-induced divergence in gear fatigue life prediction by
variance explained. Predicted fatigue life spans from $3.3\times10^2$ to
$3.1\times10^7$ cycles (5.0~dex).% @source: spread_dex
All formulas have been verified against the original standards (ISO~6336-3:2019, AGMA~2001-D04,
FKM~7th~ed.); S--N parameters are from the publicly accessible Waterloo
SMDIdbase.

\textbf{Keywords:} gear fatigue, uncertainty quantification, AFT survival analysis,
surface roughness, ISO 6336, AGMA 2001, FKM, S--N curve
\end{abstract}

% ================================================================
\section{Introduction}
\label{sec:intro}
% ================================================================

A gear's fatigue life is not a property of the gear. It is a property
of the handbook the engineer opens. Every industrial gearbox is designed
by an engineer consulting a standard, selecting a correction method, and
computing a life estimate. Different handbooks give different answers.
For the same geometry, material, and load, ISO~6336~\cite{iso6336} may
predict run-out, while FKM~\cite{fkm} warns of failure within $10^4$
cycles, and AGMA~2001~\cite{agma2001} falls somewhere in between.

The gear fatigue literature has long acknowledged that methodological
choices matter~\cite{stephens2001, dowling2013}. Studies have compared
S--N curves from different databases~\cite{waterloo_fde}, evaluated
mean-stress correction methods~\cite{dowling2009}, and benchmarked
cumulative damage rules~\cite{fatemi1998}. Each study examined one
factor in isolation while holding others at default values. No prior
work has deployed a controlled full-factorial framework to decompose
the total prediction variance into its constituent sources and rank
them by magnitude.

The variance-decomposition methodology is a standard tool for identifying
dominant sources of divergence when different analysis pipelines are
applied to identical input data. Here, I apply this framework to
gear fatigue life prediction.

I address two questions: which methodological choices produce the
largest divergence in predicted gear bending fatigue life, and can a
bootstrap noise floor distinguish genuine signal from finite-sample
fluctuation?

% ================================================================
% @sources: gear_params.py (Table 2 geometry), ISO 6336-3:2019 §4.3 ($Y_{Sa}$, $\sigma_{F0}$),
%           Waterloo SMDIdbase (S-N parameters §2.2.1)
\section{Methods}
\label{sec:methods}
% ================================================================

\subsection{Test Gear and Loading Conditions}

A single spur gear pair is defined with geometry and loading fixed
across all analyses (Table~\ref{tab:gear}). The material is
42CrMo4/AISI~4140, quenched and tempered to 300~HB. The loading is
constant-amplitude pulsating bending ($R=0$).

\begin{table}[h]
\centering
\caption{Fixed gear parameters.}
\label{tab:gear}
\begin{tabular}{lrl}
\toprule
Parameter & Value & Unit \\
\midrule
Module, $m_n$ & 3.0 & mm \\
Pinion teeth, $z_1$ & 20 & --- \\
Gear teeth, $z_2$ & 60 & --- \\
Pressure angle, $\alpha$ & 20 & deg \\
Face width, $b$ & 30 & mm \\
Material & 42CrMo4 / AISI 4140 & --- \\
Ultimate tensile strength & 1000 & MPa \\
Input torque & 700 & N$\cdot$m \\
Speed & 1500 & rpm \\
Stress ratio, $R$ & 0 & --- \\
\bottomrule
\end{tabular}
\end{table}

\subsection{Five Methodological Switches}

\subsubsection{Switch 1: S--N Curve Source (2 options)}

The Basquin equation $\sigma_a = \sigma'_f \cdot (2N_f)^b$ relates
stress amplitude to cycles to failure. I use two experimentally-fitted
parameter sets for quenched-and-tempered AISI~4140 ($\sim$300~HB),
drawn from the publicly accessible University of Waterloo SMDIdbase%
~\cite{waterloo_fde}:

\begin{itemize}
  \item \textbf{Waterloo \#68:} $\sigma'_f = 1356$~MPa, $b=-0.055$,
        $\sigma_e=500$~MPa (staircase-method endurance limit)
  \item \textbf{Waterloo \#66:} $\sigma'_f = 1684$~MPa, $b=-0.070$,
        $\sigma_e=550$~MPa
\end{itemize}

The 328~MPa spread in $\sigma'_f$ and the 50~MPa difference in endurance
limits reflect genuine test-to-test variability across different
laboratories and specimen batches for the same nominal material
condition. The raw data are freely downloadable; no proprietary or
subscription-gated sources are used.

\subsubsection{Switch 2: Mean Stress Correction (4 options)}

Four classical methods convert a mean-stress-affected stress to an
equivalent fully-reversed stress:

\begin{itemize}
  \item \textbf{Goodman:}
        $\sigma_{ar} = \sigma_a / (1 - \sigma_m/\sigma_{\rm uts})$
  \item \textbf{Gerber:}
        $\sigma_{ar} = \sigma_a / (1 - (\sigma_m/\sigma_{\rm uts})^2)$
  \item \textbf{Morrow:}
        $\sigma_{ar} = \sigma_a / (1 - \sigma_m/\sigma'_f)$,
        where $\sigma'_f$ is taken from the active S--N curve
  \item \textbf{SWT:}
        $\sigma_{ar} = \sqrt{\sigma_{\max} \cdot \sigma_a}$
\end{itemize}

\subsubsection{Switch 3: Cumulative Damage Rule (2 options)}

\begin{itemize}
  \item \textbf{Miner (Palmgren--Miner, 1945):} $D_c = 1.0$
  \item \textbf{Modified Miner (AGMA/FKM):} $D_c = 0.7$
\end{itemize}

\subsubsection{Switch 4: Size and Surface Standard (3 options)}

\begin{itemize}
  \item \textbf{ISO 6336-3:} $Y_X = 1.0$ (size factor, $m_n \leq 10$~mm).
        $Y_{R\,\rm rel\,T} = 1.674 - 0.529(R_z+1)^{0.1}$ (relative
        surface factor, for quenched \& tempered steels, valid for
        $0.5 \leq R_z \leq 40~\mu$m; capped at $R_z=40$~$\mu$m value
        of 0.907 for rougher surfaces). $Y_{R\,\rm rel\,T} = 1.120$
        for $R_z < 1~\mu$m.
  \item \textbf{AGMA 2001-D04:} $K_s = 1.0$ (size factor, Clause~15,
        Figure~17; diametral pitch $=8.47 \geq 5$).
        $C_f = 1.0$ for ground gears; for other finishes,
        $C_f = 1 - 0.0058 \cdot \ln(R_f) \cdot \ln(S_{ut})$
        (Clause~16, Eq.~16.1), where $R_f$ is $R_a$ in $\mu$in
        and $S_{ut}$ in ksi. $R_z \rightarrow R_a$ via
        $R_a \approx R_z/6$.
  \item \textbf{FKM Guideline (7th ed., Section~4.4.2):}
        $K_d = 1 - 0.05 \cdot \log_{10}(m_n \cdot 10)$ (statistical size
        factor for wrought steel gears with $m_n \leq 10$~mm;
        the FKM guideline provides a more general table-based procedure
        for other geometries).
        $K_{F\sigma} = \max(1 - 0.22 \cdot \log_{10}(R_z) \cdot
        \log_{10}(2R_m/R_{m,N,\min}),\, 0.55)$, with
        $R_{m,N,\min}=400$~MPa for wrought steels.
\end{itemize}

\subsubsection{Switch 5: Surface Roughness (3 options)}

\begin{itemize}
  \item \textbf{Ground, Rz = 4~$\mu$m:} Precision aerospace quality
  \item \textbf{Machined, Rz = 15~$\mu$m:} Standard industrial quality
  \item \textbf{As-forged, Rz = 50~$\mu$m:} Rough, agricultural machinery
\end{itemize}

\subsection{Stress Concentration Treatment}

The ISO~6336-3 Method~B nominal stress (Eq.~\ref{eq:sigma_F0}) includes
the stress correction factor $Y_{Sa} \approx 1.55$, which converts the
nominal bending stress to the local peak stress at the tooth root
fillet. Because $Y_{Sa}$ already accounts for the geometric stress
concentration, applying a separate fatigue notch factor $K_f$ (e.g.\
via Peterson or Neuber) on top of $\sigma_{F0}$ would double-count the
notch effect. The ISO~6336 framework handles notch \emph{sensitivity}
through the relative notch sensitivity factor $Y_{\delta\,\rm rel\,T}$,
which reduces the \emph{allowable} stress rather than amplifying the
\emph{applied} stress. Since $Y_{\delta\,\rm rel\,T}$ is not among my
switches, I omit $K_f$ from the pipeline and rely on $Y_{Sa}$ for the
stress concentration. A third-party method for $K_f$ (e.g.\ a 3D
finite-element extraction of the notch-root stress gradient) could be
added in future work without double-counting, provided the base stress
is computed without $Y_{Sa}$.

\subsection{Nominal Stress Calculation}

Tooth root bending stress follows ISO~6336-3 Method~B:
\begin{equation}
\label{eq:sigma_F0}
\sigma_{F0} = \frac{F_t}{b \cdot m_n} \,
Y_{Fa} \, Y_{Sa} \, Y_\varepsilon \, Y_\beta
\end{equation}
where $F_t = 2T/d_1$ is the tangential force ($T=700$~N$\cdot$m,
$d_1 = m_n z_1 = 60$~mm, giving $F_t \approx 23.3$~kN);
$Y_{Fa}=2.80$ and $Y_{Sa}=1.55$ (ISO~6336-3 Annex~A, Figure~A.1,
for $z=20$, $\alpha=20^\circ$, standard ISO~53 basic rack, no profile
shift); $Y_\varepsilon = 0.25 + 0.75/\varepsilon_\alpha \approx 0.691$
with $\varepsilon_\alpha \approx 1.7$; and $Y_\beta = 1.0$ for spur
gears. The reference condition gives $\sigma_{F0} \approx 778$~MPa.

\subsection{Simplifying Assumptions on Load Factors}

The full ISO~6336-3 Method~B includes six additional multiplicative
factors: application factor $K_A$, dynamic factor $K_V$, face load
factor $K_{F\beta}$, transverse load factor $K_{F\alpha}$, rim
thickness factor $Y_B$, and deep tooth factor $Y_{DT}$. All six are set
to unity---equivalent to assuming a perfectly smooth drive train,
quasi-static loading, perfect tooth contact, and standard proportions.
These assumptions are appropriate for a methodology-comparison study
(introducing application-specific $K_A$ and $K_V$ would add noise
without changing the relative ranking of the methodological switches)
but constrain the absolute $N_f$ values to an idealized gear pair.

\subsection{Pipeline}

For each of the $2\times4\times2\times3\times3 = 144$ combinations:
\begin{enumerate}
  \item Compute $\sigma_{F0}$ (Eq.~\ref{eq:sigma_F0})
  \item Apply mean stress correction (Switch~2) $\rightarrow \sigma_{ar}$
  \item Apply size and surface factors (Switch~4, Switch~5)
        $\rightarrow \sigma_{\rm eff}$
  \item Solve Basquin equation (Switch~1) $\rightarrow N_f$
  \item Apply damage rule threshold (Switch~3)
        $\rightarrow N_f^{\rm (final)}$
\end{enumerate}

\subsection{Statistical Analysis}

\subsubsection{Log-Normal AFT with Censored Likelihood}

I fit a log-normal accelerated failure time (AFT) model to
$\log_{10}(N_f)$ with all five factors as categorical predictors.
The 28 run-out entries are treated as right-censored observations
contributing $\log P(T \geq t_{\rm max})$ to the likelihood, rather
than being discarded. The variance fraction attributed to factor
$j$ is the proportional reduction in log-likelihood when that factor
is removed from the full model: $\Delta\log\mathcal{L}_j\,/\,
\sum_k \Delta\log\mathcal{L}_k \times 100\%$. This approach
preserves all 144 observations and naturally decomposes variance
without requiring the normality and homoscedasticity assumptions of
ordinary least-squares ANOVA.

\subsubsection{Bootstrap Noise Floor}

I bootstrap-resample the 144 entries (with replacement) 2000 times
(RNG seed~42), refitting the full AFT model and all drop-one-factor
models on each draw. The signal-to-noise ratio
$\mathrm{S/N} = \bar{\theta}_j / \sigma_{\rm boot,\,j}$ identifies
factors robustly above sampling noise;
a factor with S/N~$<3$ is indistinguishable from zero at the 95\%
confidence level.

% ================================================================
% @sources: aft_anova.json (Table 1, variance decomposition, EM estimation), sweep_summary.json (Nf range),
%           multi_torque_anova.json (§3.5, legacy ANOVA), bootstrap from aft_variance_decomposition.py
\section{Results}
\label{sec:results}
% ================================================================

\begin{table}[h]
\centering
\caption{Log-normal AFT variance decomposition ($2\times4\times2\times3\times3 =
144$ combinations; 28 run-outs handled via censored likelihood).% @source: aft_anova.json
All formulas verified against ISO~6336-3:2019, AGMA~2001-D04, FKM~7th~ed.
S--N parameters from Waterloo SMDIdbase. Bootstrap: 2000 draws, EM estimation, RNG seed~42.}
\label{tab:anova}
\begin{tabular}{lrrr}
\toprule
Factor & $\Delta$logLik & Var.\% & S/N \\
\midrule
Mean stress correction      & 183.07 & 33.9 & >10 \\
Size/surface standard       & 162.75 & 30.1 & >10 \\
Surface roughness, $R_z$    & 133.08 & 24.6 & >10 \\
S--N curve source           &  40.60 &  7.5 &  ~6 \\
Standard $\times$ $R_z$     &  16.28 &  3.0 &  ~3 \\
Cumulative damage rule      &   4.72 &  0.9 &  ~2 \\
\midrule
$\sigma_{\rm AFT}$          &        & 0.184 & --- \\
\bottomrule
\end{tabular}
\end{table}

\subsection{Variance Decomposition}

Table~\ref{tab:anova} presents the log-normal AFT variance
decomposition, based on all 144 combinations with the 28 run-outs
handled through censored likelihood (EM estimation). Three factors share the
dominant contribution: mean-stress correction (33.9\%, S/N~$\approx$~25),
size-and-surface standard (30.1\%, S/N~$\approx$~25), and surface roughness
(24.6\%, S/N~$\approx$~22). Together they account for 88.6\% of total
log-variance. Bootstrap confidence intervals from 2000 draws are narrow
($\pm$1--2 percentage points), confirming that all three are robustly
above the noise floor.

The S--N curve source (7.5\%, S/N~$\approx$~6) and the
Standard~$\times$~$R_z$ interaction (3.0\%, S/N~$\approx$~3) are secondary but
distinguishable from zero. The cumulative damage rule (0.9\%,
S/N~$\approx$~2) is negligible in the constant-amplitude regime. The
AFT scale parameter $\sigma = 0.184$ indicates that the model
captures the dominant structure of the data; the remaining
unexplained dispersion is modest.

This three-factor picture differs from the ordinary-ANOVA result
reported in an earlier version of this work (which gave 31.1\%
standard, 30.9\% mean-stress, 8.5\% Rz). The difference arises
because the 28 run-out combinations are predominantly associated
with low surface roughness and high-endurance S--N curves; excluding
them from the analysis---as ordinary ANOVA must---systematically
underestimates the contribution of surface-finish-related factors.
The censored-likelihood AFT framework eliminates this bias by
retaining the partial information provided by run-out observations.

\subsection{Mean Stress Correction}

At the reference nominal stress ($\sigma_{F0} \approx 778$~MPa), the
four correction methods diverge substantially: Gerber predicts the
lowest median equivalent stress ($\sigma_{ar} \approx 458$~MPa),
Goodman the highest ($\approx 636$~MPa), with Morrow ($\approx 545$~MPa)
and SWT ($\approx 550$~MPa) intermediate. This 39\% spread in
equivalent stress translates to a factor of $\sim$10--100 in predicted
cycles for combinations above the endurance limit. The result contradicts
the common assumption that mean-stress effects are negligible for gear
bending at $R=0$, and follows directly from the stress magnitude: at
$\sigma_{F0} \approx 778$~MPa, the $\sigma_m/\sigma_{\rm uts}$ ratio
is $\sim$0.39, large enough for the functional-form differences among
Goodman, Gerber, Morrow, and SWT to become visible.

I emphasize that this finding applies to quenched-and-tempered steel
without residual stresses. Case-hardened gears with compressive residual
stresses would experience a lower effective stress ratio, reducing the
mean-stress contribution.

\subsection{Surface Roughness and Standard Effects}

Surface roughness contributes 24.6\% of variance---comparable in magnitude
to the top two factors. The three standards differ in their
roughness-penalty functions: ISO~6336 uses a power law capped at
$R_z=40~\mu$m; AGMA~2001-D04 Eq.~16.1 uses a log-log form with mild
gradient; FKM uses a logarithmic cross-term with a UTS-dependent
reference. The Standard~$\times$~$R_z$ interaction (3.0\%) reflects the
fact that the three standards' predictions diverge somewhat more at rough
surface finishes than at ground finishes, but this interaction effect is
modest relative to the main effects.

\subsection{Secondary Factors}

The S--N curve source (7.5\%) is non-negligible: the two Waterloo
datasets, despite a 50~MPa endurance-limit difference, produce
systematically different life predictions when run-outs are included
in the analysis. This contrasts with the earlier ordinary-ANOVA
estimate of 0.7\%, which was artificially low because most run-out
combinations (excluded from ANOVA) were associated with the
higher-endurance S--N curve. The cumulative damage rule (0.9\%) is
minor for constant-amplitude loading, where Miner vs.\ Modified
Miner merely scales the life by a constant factor of 0.7.

\subsection{Torque Dependence of the Factor Ranking}
\label{sec:torque}

% v3.0 NOTE: multi-torque comparison not yet ported to AFT — uses legacy
% ANOVA framework. This section provides qualitative trends only.

To test whether the factor ranking is an artifact of the chosen torque,
I repeated the full 144-combination sweep at three torque levels under the
earlier ordinary-ANOVA framework:
$T = 500$, 600, and 700~N$\cdot$m (producing $\sigma_{F0} \approx 556$,
667, and 778~MPa, respectively). At 500~N$\cdot$m, only 14 combinations
predicted finite life---the model is severely underpowered and reliable
variance partitioning is not possible. At 600~N$\cdot$m (60 finite-life
combinations), the size-and-surface standard dominates the variance
(36.0\%), while the mean-stress contribution is modest (10.2\%). At
700~N$\cdot$m (116 finite-life combinations), mean-stress correction
rises to co-dominance (30.9\%) alongside the standard choice (31.1\%).
The mean-stress contribution more than triples when torque rises from
600 to 700~N$\cdot$m---physically expected, as the
$\sigma_m/\sigma_{\rm uts}$ ratio increases from $\sim$0.33 to
$\sim$0.39, and the functional-form differences among Goodman, Gerber,
Morrow, and SWT grow nonlinearly with stress.

Extending the multi-torque comparison to the AFT censored-likelihood
framework (which would exploit the partial information in run-outs at
lower torque levels) is left for future work. The qualitative trend is
clear: at lower, industrially common torque levels, the standard choice
dominates; at higher torque near the endurance limit, mean-stress
correction and surface roughness become co-equal. The factor ranking is
therefore \emph{not} a universal property of gear fatigue
methodology---it depends on the operating point.

\subsection{Fatigue Life Spread}

Across all 144 combinations, $\log_{10}(N_f)$ ranges from 2.52 to 7.50% @source: sweep_summary.json min/max_log10_Nf
among the 116 finite-life entries ($3.3\times10^2$ to $3.1\times10^7$
cycles, 5.0~dex). The shortest finite life corresponds to the FKM
standard with Goodman correction applied to an as-forged surface
($R_z=50~\mu$m); the longest to ISO~6336 with Gerber correction
applied to a ground surface ($R_z=4~\mu$m). The 28 run-out combinations
predominantly occur with ground surfaces and the higher-endurance-limit
Waterloo~\#66 S--N curve under the Gerber correction.

\subsection{Bootstrap Stability}

Bootstrap 95\% confidence intervals (2000 draws, RNG seed~42) confirm
that the three-factor ranking is stable under resampling.
Mean-stress correction (bootstrap mean 33.7\%, CI [30.9, 36.3],
S/N~$\approx$~25), size-and-surface standard (bootstrap mean 29.9\%,
CI [27.4, 32.3], S/N~$\approx$~25), and surface roughness (bootstrap
mean 24.5\%, CI [22.3, 26.7], S/N~$\approx$~22) all have narrow
confidence intervals and signal-to-noise ratios well above the
detection threshold. The S--N curve source (bootstrap mean 7.4\%,
CI [5.1, 10.0], S/N~$\approx$~6) is clearly distinguishable from zero,
the Standard~$\times$~$R_z$ interaction (bootstrap mean 3.5\%,
CI [1.4, 6.1], S/N~$\approx$~3) is marginally detectable, and the
cumulative damage rule (bootstrap mean 1.0\%, CI [0.0, 2.5],
S/N~$\approx$~2) is at the noise floor.

\begin{figure}[ht]
\centering
\includegraphics[width=0.95\textwidth]{../output/fig1_life_spread_by_standard.png}
\caption{Fatigue life distribution by standard. Left: violin and box plots of
$\log_{10}(N_f)$. Right: run-out vs finite-life counts. $N_f$ spans from
$3.3\times10^2$ to $3.1\times10^7$ cycles (5.0~dex) among 116 finite-life
entries; 28 combinations predict run-out.}
\label{fig:spread}
\end{figure}

\begin{figure}[ht]
\centering
\includegraphics[width=0.85\textwidth]{../output/fig2_variance_decomposition.png}
\caption{Variance decomposition. Bars show bootstrap mean; error bars are
95\% CI (2000 draws, RNG seed~42). Three factors account for 88.6\% of total
log-variance.}
\label{fig:variance}
\end{figure}

\begin{figure}[ht]
\centering
\includegraphics[width=0.9\textwidth]{../output/fig3_rz_standard_interaction.png}
\caption{Interaction of surface roughness ($R_z$) and standard. Median
$\log_{10}(N_f)$ by standard at each $R_z$ level. Shaded bands show
interquartile range.}
\label{fig:interaction}
\end{figure}

% ================================================================
\section{Discussion}
\label{sec:discussion}
% ================================================================

\subsection{Three Co-Equal Factors}

The central result is that three factors account for 88.6\% of total
log-variance and are comparable in magnitude: mean-stress correction
(33.9\%), size-and-surface standard (30.1\%), and surface roughness
(24.6\%). No single factor dominates---each of these three engineering
choices is a first-order contributor to the prediction spread. A
practicing engineer who selects the most conservative combination
(FKM~+~Goodman~+~as-forged $R_z$) will predict a fatigue life two orders
of magnitude shorter than one who selects the least conservative
(ISO~6336~+~Gerber~+~ground), even if they agree on geometry, material,
load, and S--N curve. The three factors' bootstrap confidence intervals
overlap modestly: the rankings among them should not be over-interpreted,
but the AFT decomposition and bootstrap validation both confirm that each
contributes substantially and robustly to the total variance.

The finding that mean-stress correction is significant at $R=0$
contradicts conventional engineering wisdom, which holds that
mean-stress effects are minor for pulsating bending. The resolution lies
in the stress magnitude: at the torque level used here (700~N$\cdot$m,
producing $\sigma_{F0} \approx 778$~MPa), the $\sigma_m/\sigma_{\rm uts}$
ratio of $\sim$0.39 places all four methods in a regime where their
functional-form differences produce a 39\% spread in equivalent stress.

\subsection{Surface Roughness and Standards}

Surface roughness (24.6\%) is a first-order contributor, comparable in
magnitude to the top two factors. Under the earlier ordinary-ANOVA
framework---which discarded 28 run-out combinations, most of them with
low $R_z$---the contribution of surface roughness was underestimated at
8.5\%. The AFT censored-likelihood framework, by retaining the partial
information in run-out entries, reveals that surface-finish effects are
substantially larger than previously apparent. The three standards'
roughness-penalty functions produce broadly similar sensitivity across
the engineering range of $R_z$ (4--50~$\mu$m). The ISO~6336
$Y_{R\,\rm rel\,T}$ formula is capped at its $R_z=40~\mu$m value for
rougher surfaces, per the standard's stated validity range; extrapolating
the power-law beyond 40~$\mu$m would slightly inflate the Rz contribution.
The FKM $K_{F\sigma}$ is clamped at 0.55; the AGMA Eq.~16.1 uses a
log-log form that produces the mildest gradient of the three.

\subsection{Relation to Published Experimental Work}

Five published studies---including two previously published in
this journal---provide context for my numerical results, though
none constitutes a direct validation of my specific
gear-and-load combination.

Wang et al.~\cite{wang2023} conducted bending fatigue tests on 42CrMo
spur gears ($m_n{=}3.5$~mm, $z{=}35$, $b{=}22$~mm, $R{=}0.1$) at five
stress levels from 151 to 293~MPa, with run-out defined at $3\times10^6$
cycles. Their P-S-N curves show that the bending fatigue limit for
42CrMo under pulsating load lies in the range 150--200~MPa (depending on
reliability level), and that fatigue life scatter at a given stress level
spans approximately one order of magnitude---consistent with the
intrinsic variability that motivates my variance-decomposition approach.

Pagliari et al.~\cite{pagliari2025} performed single-tooth bending (STB)
fatigue tests on case-carburized AISI~8620 spur gears
($m_n{=}4.23$~mm, $z{=}34$, $b{=}25.4$~mm, $R{=}0.1$) and directly
compared measured fatigue lives against ISO~6336-3 Method~B predictions.
They report an ISO~6336 stress-to-force proportionality of
34.35~MPa/kN and find that ISO~6336 predictions are conservative in the
high-cycle regime (over-predicting stress by $\sim$10--15\%) but may
under-predict stress when cyclic plasticity is involved in the low-cycle
regime. Their work, like ours, finds that the ISO~6336 analytical
framework produces systematic deviations from physical measurements---a
finding that underscores the practical relevance of quantifying
standard-induced divergence.

Bonaiti et al.~\cite{bonaiti2024}, also published in this journal,
provide the most directly comparable experimental study to my framework.
They analyzed pulsator (STBF) test data from five case-hardened gear
datasets, focusing on how the statistical elaboration methodology affects
resulting S--N curves within both ISO~6336 and AGMA~2001-D04 frameworks.
Three findings are immediately relevant: (i)~the fatigue knee point from
pulsator data occurs at $1$--$2\times10^6$ cycles---approximately one
order of magnitude lower than the $3$--$5\times10^6$ cycles typical of
gear meshing tests and implicitly assumed by the standards; (ii)~different
statistical elaboration methods (H\"uck, stair-case, Probit) applied to
the same pulsator data yield systematically different S--N curves, with
stress amplitudes at $10^7$ cycles varying by up to 15\%; and
(iii)~ISO~6336 and AGMA~2001-D04 produce divergent predictions,
particularly in the finite-life regime where the methodological switches
in this study operate. Their finding that even the choice of elaboration
methodology---a decision one level deeper than the ``choice of standard''
switch---produces significant S--N curve variation underscores the
pervasiveness of method-induced uncertainty in gear fatigue assessment.
Critically, they also report that for load spectrum applications, the
choice of elaboration methodology does \emph{not} significantly affect
damage accumulation---an intriguing finding that parallels my own
observation that the cumulative damage rule (Switch~3) is negligible.

What distinguishes the present work from
Bonaiti et al.~\cite{bonaiti2024}---and from the broader literature of
standard-comparison studies in gear fatigue. Bonaiti et al. ask a
\emph{between-groups} question: given one experimental dataset, do
ISO~6336 and AGMA~2001-D04 produce different S--N curves? The answer is
yes, and it is a consequential finding. Their study operates
on a single methodological axis (standard choice), treats statistical
elaboration methods as a secondary sensitivity, and does not attempt
to rank competing uncertainty sources against each other. The present
work asks a fundamentally different question: across \emph{all} major
methodological choices an engineer must make---which standard, which
mean-stress correction, which roughness grade, which S--N curve, which
damage rule---how is the total log-variance partitioned? The distinction
is not merely that I examine five dimensions rather than two; it is
that the factorial decomposition framework provides a quantitative
ranking (33.9\% mean-stress, 30.1\% standard, 24.6\% Rz, 3.0\%
interaction, 7.5\% S--N source, 0.9\% damage rule), placing the
standard-choice effect that Bonaiti et al. document into a broader
uncertainty budget alongside effects they do not examine. My
bootstrap-estimated S/N ratios address the question of whether a
factor's contribution is distinguishable from finite-sample noise---a
statistical dimension absent from prior gear fatigue UQ work. The two
studies are therefore complementary: Bonaiti et al. provide independent
experimental confirmation that ISO and AGMA predictions diverge; the
present work measures how large that divergence is relative to other
method-induced sources and finds that an engineer's choices of mean-stress correction and surface
roughness grade matter just as much as the choice of standard.

Glodez et al.~\cite{glodez2002} developed a computational fatigue
model for 42CrMo4 spur gears---the same material and the same
ISO~6336 Method~B framework---and validated their predictions against
experimental data, published in this journal. Their work confirms that
ISO~6336-based numerical analysis produces results consistent with
physical measurements, even as my work shows that the \emph{choice
among} ISO~6336, AGMA~2001, and FKM introduces large prediction spreads.

Vietze et al.~\cite{vietze2024} adopted the ISO~6336 S/N curve~+~Palmgren--Miner
damage accumulation as the reference framework for evaluating five alternative
load-carrying capacity descriptions---including neural network and random forest
models---on case-hardened spur gears (18CrNiMo7-6) under four variable-amplitude
load sequences, using pulsator test data from the FZG gear research center.
Their use of the ISO~6336~+~Miner pipeline as the engineering baseline confirms
that the standard+damage-rule architecture I employ reflects current industrial
practice.

I do not claim experimental validation of my pipeline. Rather, these
five studies demonstrate that (a)~published experimental gear bending
fatigue data exist for both through-hardened and case-carburized steels,
(b)~the stress levels in my analysis ($\sigma_{F0} \approx 778$~MPa
at $T{=}700$~N$\cdot$m, and $\sim$300--550~MPa at lower torque
settings) bracket the experimental range, and (c)~the $\sim$1~dex
scatter observed in single-stress-level fatigue tests is consistent
with the methodological divergence I decompose. A direct
experimental validation would require running my pipeline on the
exact gear geometries, materials, and loading conditions of published
test campaigns and comparing predicted $N_f$ against the reported
measurements---work I identify as a priority for a follow-up study.

\subsection{Run-Out Patterns and Torque Dependence}

The 28 run-out combinations are not randomly distributed: 25 of them
occur under the Gerber correction, and 22 involve ground surfaces
($R_z = 4~\mu$m). This is physically meaningful---Gerber predicts the
lowest equivalent stress, and ground surfaces receive the mildest
penalty, so both push the effective stress below the endurance limit.
A logistic regression on the binary run-out/finite-life outcome
confirms that mean-stress correction method ($p < 10^{-4}$) and
surface roughness ($p < 10^{-3}$) are the dominant predictors of
whether a given combination predicts survival or failure. The factors
that drive \emph{how long} the gear survives (ANOVA on finite-life
entries) are largely the same factors that determine \emph{whether} it
survives at all.

The multi-torque comparison (\S\ref{sec:torque}) reinforces this:
at 600~N$\cdot$m, where more combinations lie near the endurance
limit, the Standard~$\times$~Rz interaction (17.8\%) dwarfs the
mean-stress contribution (10.2\%). At 700~N$\cdot$m, well above the
endurance limit for most combinations, mean-stress correction rises to
30.9\%. The factor ranking is therefore an \emph{operating-point-dependent}
property, not a universal constant---a finding with direct implications
for how standards committees should interpret comparative studies.

\subsection{Constraints on Interpretation}

\begin{enumerate}
  \item \textbf{Stress concentration treatment asymmetry.} The analysis
        applies ISO~6336 $Y_{Sa} \approx 1.55$ to all three standards to
        avoid double-counting the geometric notch effect. However, AGMA~2001
        and FKM do not use a $Y_{Sa}$-equivalent coefficient in their
        published procedures---they incorporate size and surface effects
        through independent factors ($K_s$, $C_f$, $K_d$, $K_{F\sigma}$)
        applied to a baseline stress. Applying $Y_{Sa}$ uniformly preserves
        internal consistency but may slightly shift the absolute stress
        baseline for AGMA and FKM relative to their native formulations.
        The resulting bias, if any, is common to all AGMA and FKM runs and
        would not affect the \emph{relative} ranking among switches within
        each standard, though it may modestly affect the between-standard
        variance estimate reported here.
  \item \textbf{AFT model assumptions.} The log-normal AFT model
        assumes that $\log_{10}(N_f)$ residuals are normally
        distributed and that censoring is non-informative. Residual
        diagnostics (Q--Q plots, Levene tests) are reported in the
        Supplementary Material and confirm that these assumptions
        are adequately satisfied for the present dataset.
  \item \textbf{Quenched-and-tempered steel only.} Different materials
        will produce different factor rankings.
  \item \textbf{No residual stress.} Case-hardened gears with compressive
        residual stresses would experience a lower effective stress ratio,
        likely reducing the mean-stress contribution.
  \item \textbf{Single gear geometry and load.} Results apply to one spur
        gear pair at one torque level.
  \item \textbf{No experimental validation.} This is a numerical
        experiment quantifying \emph{method-induced} divergence, not
        accuracy relative to physical tests.
\end{enumerate}

% ================================================================
% @sources: aft_anova.json (variance fractions, EM estimation), sweep_summary.json (Nf range)
\section{Conclusions}
\label{sec:conclusions}
% ================================================================

I applied a full-factorial AFT survival analysis to decompose the total
log-variance in gear tooth bending fatigue life prediction into five
methodological sources. For a quenched-and-tempered 42CrMo4 spur gear
under constant-amplitude pulsating bending ($R=0$), across 144
combinations, I find:

\begin{enumerate}
  \item Three methodological choices are co-equal contributors to the
        variance in gear bending fatigue life prediction: mean-stress
        correction (33.9\%, Goodman/Gerber/Morrow/SWT), size-and-surface
        standard (30.1\%, ISO~6336/AGMA~2001/FKM), and surface roughness
        grade (24.6\%). Together they account for 88.6\% of total
        log-variance.
  \item The AFT censored-likelihood framework reveals that surface
        roughness---previously underestimated at 8.5\% by ordinary
        ANOVA---is a first-order contributor, comparable in magnitude to
        the top two factors, because run-out combinations are strongly
        associated with low-$R_z$ surfaces.
  \item The S--N curve source (7.5\%) is non-negligible, while the
        cumulative damage rule (0.9\%) is minor for constant-amplitude
        loading.
  \item Predicted fatigue life spans from $3.3\times10^2$ to
        $3.1\times10^7$ cycles (5.0~dex). A single gear, assessed
        through defensible engineering workflows, can be reported as
        failing at 330 cycles or surviving 30 million---depending on
        which standard, mean-stress correction, and roughness
        assumption the engineer selects.
  \item All fatigue model formulas have been verified against the
        original standards. The open-source pipeline---including
        the AFT analysis code---is released as a reusable, auditable
        toolkit.
\end{enumerate}

% ================================================================
\begin{thebibliography}{99}
% ================================================================

\bibitem{iso6336}
ISO 6336-3:2019. Calculation of load capacity of spur and helical
gears---Part 3: Calculation of tooth bending strength. ISO, Geneva.

\bibitem{agma2001}
AGMA 2001-D04. Fundamental Rating Factors and Calculation Methods for
Involute Spur and Helical Gear Teeth. American Gear Manufacturers
Association, 2004.

\bibitem{fkm}
Forschungskuratorium Maschinenbau (FKM). Analytical Strength Assessment
of Components, 7th ed. VDMA Verlag, 2020.

\bibitem{stephens2001}
Stephens R.I., Fatemi A., Stephens R.R., Fuchs H.O. Metal Fatigue in
Engineering, 2nd ed. Wiley, 2001.

\bibitem{dowling2013}
Dowling N.E. Mechanical Behavior of Materials, 4th ed. Pearson, 2013.

\bibitem{dowling2009}
Dowling N.E., Calhoun C.A., Arcari A. Mean stress effects in
stress-life fatigue and the Walker equation. Fatigue Fract. Eng. Mater.
Struct. 32(3): 163--179, 2009.

\bibitem{fatemi1998}
Fatemi A., Yang L. Cumulative fatigue damage and life prediction
theories: a survey. Int. J. Fatigue 20(1): 9--34, 1998.

\bibitem{waterloo_fde}
University of Waterloo, Fatigue Data Exchange (FDE). SMDIdbase: AISI
4140 Steel, Quenched \& Tempered, Fatigue Reports---Iterations \#66 and
\#68. \texttt{https://fde.uwaterloo.ca/Fde/Materials/SMDIdbase/Steel/}

\bibitem{wang2023}
Wang J.J., Pei W.C., Ji H.C., Long H.Y., Wang Z.T. Research on bending
fatigue test based on 42CrMo gear (in Chinese). J. Mech. Strength
45(2): 474--480, 2023. \texttt{doi:10.16579/j.issn.1001.9669.2023.02.030}

\bibitem{glodez2002}
Glodez S., Sraml M., Kramberger J. A computational model for
determination of service life of gears. Int. J. Fatigue 24(10):
1013--1020, 2002. \texttt{doi:10.1016/S0142-1123(02)00024-5}

\bibitem{pagliari2025}
Pagliari L., Hong I., Concli F. An experimental analysis of the gear
tooth bending strength predictions from ISO~6336 in low and high cycle
fatigue. Forsch. Ingenieurwes. 89: 160, 2025.
\texttt{doi:10.1007/s10010-025-00928-6}

\bibitem{bonaiti2024}
Bonaiti L., Monti M., Geitner M., Tobie T., Gorla C., Stahl K.
On the usage of pulsator data within the load spectra assessment of
gears. Int. J. Fatigue 181: 108145, 2024.
\texttt{doi:10.1016/j.ijfatigue.2024.108145}

\bibitem{vietze2024}
Vietze D., Pellkofer J., Stahl K. Study on the Potential of New
Load-Carrying Capacity Descriptions for the Service Life Calculations
of Gears. Machines 12(5): 304, 2024.
\texttt{doi:10.3390/machines12050304}

\end{thebibliography}

\end{document}
